% Chapter 18: The Cap Set Problem
\let\LocalMacro\relax

\chapter{Cap Set Problem}

\section{The Problem}

\subsection{Informal}
Imagine a big grid of points in space where the arithmetic wraps around like a clock. You want to pick as many points as possible, but with one strict rule: no three of your chosen points can be equally spaced in a line. How many points can you pick before you are forced to have three in a row? This problem stumped mathematicians for decades, despite being easy to state.

\subsection{Formal}
\begin{conjecture}[Frankl-Graham-Rödl, 1987]
$r_3(n) = o(3^n)$.
\end{conjecture}

The trivial bound is $3^n$. Meshulam (1995) proved $r_3(n) \le 2 \cdot 3^n / n$, using Fourier analysis. The best pre-2016 bound was $r_3(n) \le 3^n / n^{1+c}$ for a small constant $c > 0$.

\subsection{Historical Context}
The cap set problem is the $\mathbb{F}_3$-analogue of Roth's theorem on arithmetic progressions in the integers (1953). While Roth's theorem was proved using Fourier analysis on $\mathbb{Z}$, the finite-field version resisted analogous techniques for decades.

\subsection{What Was Known}
Roth proved in 1953 that subsets of the integers without three-term arithmetic progressions must be thin. Meshulam extended this to three-dimensional grid arithmetic in 1995, proving that cap sets must be significantly smaller than the full grid. But the best known bound was still far from the conjectured answer. The polynomial method, newly adapted to combinatorics by Croot, Lev, and Pach in 2016, provided the breakthrough.

\subsection{Why It Matters}
The cap set problem is a testing ground for new mathematical techniques. The slice rank method used to solve it has since been applied to many other problems in combinatorics and theoretical computer science. Understanding how large structured sets can be has applications to coding theory and the study of error-correcting codes.

\subsection{Simplified Version}
The classical version of this problem works over the integers: Roth's theorem (1953) says that any subset of the integers $\{1, 2, \ldots, N\}$ without three-term arithmetic progressions must be thin---its size is at most $N / \log \log N$. The cap set problem asks the same question over $\mathbb{F}_3^n$ (the $n$-dimensional vector space over the field with 3 elements), where the arithmetic wraps around like a clock. The polynomial method, adapted from finite-field geometry, finally cracked this harder version.

\section{Mathematical Theory}

\subsection{Prerequisites}
This chapter assumes familiarity with:
\begin{itemize}
\item Basic combinatorics (binomial coefficients, extremal combinatorics)
\item Linear algebra (rank of matrices and tensors)
\item Finite fields ($\mathbb{F}_3$, arithmetic in characteristic 3)
\end{itemize}

\subsection{The Polynomial Method}

The polynomial method is a versatile combinatorial tool with a simple core idea: to show that a set $A$ is small, find a nonzero polynomial that vanishes on all of $A$. If the polynomial has low degree, it cannot have too many zeros, so $A$ must be small.

\begin{definition}[Vanishing Polynomial]
A polynomial $P \in \mathbb{F}_q[x_1, \ldots, x_n]$ \emph{vanishes} on a set $A \subset \mathbb{F}_q^n$ if $P(a) = 0$ for all $a \in A$.
\end{definition}

The key fact connecting polynomials to finite fields is:

\begin{fact}[Schwartz-Zippel]
A nonzero polynomial of degree $d$ over $\mathbb{F}_q^n$ has at most $d \cdot q^{n-1}$ zeros.
\end{fact}

Dvir's 2008 proof of the Kakeya conjecture over finite fields demonstrated the power of this approach: by constructing a low-degree polynomial vanishing on a Kakeya set, he showed such sets must be large. The cap set problem asks whether a similar strategy can beat Meshulam's Fourier-analytic bound.

\subsection{Tensor Rank and Slice Rank}

The breakthrough of Croot, Lev, Pach, and Ellenberg--Gijswijt reformulates the cap set problem in terms of \emph{tensor rank}, a generalization of matrix rank.

\begin{definition}[Tensor]
A 3-tensor $T \in V_1 \otimes V_2 \otimes V_3$ over a field $\mathbb{F}$ is a multilinear map $T: V_1^* \times V_2^* \times V_3^* \to \mathbb{F}$.
\end{definition}

\begin{definition}[Slice Rank]
The \emph{slice rank} of a 3-tensor $T \in V_1 \otimes V_2 \otimes V_3$ is the minimum $r$ such that $T$ can be written as a sum of $r$ terms, each of which is a product of a vector in one factor with a bilinear form on the other two:
\[
T = \sum_{i=1}^{r} v_i \otimes B_i
\]
where each $v_i$ lives in one of $V_1, V_2, V_3$ and $B_i$ lives in the tensor product of the other two.
\end{definition}

Slice rank is always at most the tensor rank, and for diagonal tensors (those supported only on the ``diagonal'' $x = y = z$), slice rank equals the size of the diagonal support.

\subsection{The Cap Set Tensor}

For the cap set problem over $\mathbb{F}_3^n$, define the tensor:
\[
T(x,y,z) = \begin{cases} 1 & \text{if } x + y + z = 0 \\ 0 & \text{otherwise} \end{cases}
\]
A subset $A \subset \mathbb{F}_3^n$ is a cap set if and only if $T|_{A \times A \times A}$ is diagonal (i.e., $T(x,y,z) \neq 0$ only when $x = y = z$). The slice rank of $T$ then upper-bounds $|A|$:
\[
|A| \le \text{slice rank}(T).
\]

\begin{highlightbox}
The key insight: the cap set condition is exactly the condition that the restriction of $T$ to $A$ is diagonal. Since diagonal tensors have slice rank equal to the size of the diagonal, any cap set has size at most the slice rank of $T$.
\end{highlightbox}

\subsection{Counting Monomials}

The slice rank of $T$ is bounded by counting low-degree monomials. Consider the vector space $\mathcal{P}_{<2n/3}$ of polynomials in $\mathbb{F}_3[x_1, \ldots, x_n]$ of degree less than $2n/3$ (with each variable exponent in $\{0, 1, 2\}$). Using Stirling's approximation:

\begin{theorem}[Monomial Count]
$|\mathcal{P}_{<2n/3}| \le 2.756^n$.
\end{theorem}

The slice rank of $T$ is at most $3 \cdot |\mathcal{P}_{<2n/3}|$ (the factor 3 accounts for the three tensor factors), giving:
\[
r_3(n) \le 3 \cdot 2.756^n.
\]

\section{The Solution}

\subsection{What Was Missing: Beating the Fourier Barrier}

The cap set problem asks: how large can a subset $A \subset \mathbb{F}_3^n$ be if it contains no three-term arithmetic progression (i.e., no three distinct points $x, y, z$ with $x + y + z = 0$)?

The best result before 2016 was Meshulam's 1995 bound:
\[
r_3(n) \le \frac{2 \cdot 3^n}{n},
\]
obtained via Fourier analysis on $\mathbb{F}_3^n$. Incremental improvements nudged this to $3^n / n^{1+c}$ for a small constant $c > 0$, but nobody could break past the $2.756^n$ barrier. The gap was clear: Fourier analysis could remove a polynomial factor from $3^n$, but the conjecture (due to Frankl, Graham, and R\"odl) predicted $r_3(n) = o(3^n)$---indeed, the truth should be \emph{subexponential} in $n$, meaning the ratio $r_3(n)/3^n$ should tend to zero faster than any exponential.

\subsection{Why Existing Methods Failed}

Fourier analysis works beautifully for additive combinatorics over $\mathbb{Z}$ (Roth's 1953 theorem) and has been the workhorse for decades. The obstacle in $\mathbb{F}_3^n$ is structural:

\begin{itemize}
\item \textbf{The tensor structure of cap sets}: The condition ``no three points in arithmetic progression'' is a \emph{three-body} constraint: it involves triples $(x, y, z)$ with $x + y + z = 0$. Fourier analysis handles two-body correlations (via the Fourier transform of a single function) but cannot directly control three-body correlations without losing a factor of $3^n$.

\item \textbf{Pl\"unnecke-Ruzsa and entropy methods}: These tools, powerful for sumset problems, gave bounds of the form $|A| \le C \cdot 3^n / n^c$ for some small $c$. The problem is that these methods are ``lossy''---each step of the argument discards information about the higher-order structure of $A$.

\item \textbf{The polynomial method was untested}: Before 2016, the polynomial method (vanishing polynomials, Schwartz-Zippel) had been used mainly in geometry (Dvir's 2008 Kakeya proof) and coding theory. Nobody had systematically applied it to additive combinatorics problems over $\mathbb{F}_3^n$.
\end{itemize}

The fundamental obstacle was that cap sets live in a world where \emph{the interesting structure is three-dimensional} (triples summing to zero), but all available tools were fundamentally two-dimensional (Fourier analysis, entropy, sumset inequalities).

\subsection{What Was Novel: The Slice Rank Method}

The breakthrough came in two stages in 2016:

\begin{enumerate}
\item \textbf{Croot, Lev, and Pach} (June 2016) used the polynomial method to prove that progression-free subsets of $\mathbb{Z}_4^n$ are exponentially small---a breakthrough for a related but easier problem.

\item \textbf{Ellenberg and Gijswijt} (July 2016) adapted and generalized the idea to $\mathbb{F}_3^n$ using \emph{slice rank}, a notion of rank for tensors.
\end{enumerate}

The key innovation is a three-step reduction:

\textbf{Step 1: From cap sets to diagonal tensors.} Define the tensor $T \in \mathbb{F}_3^n \otimes \mathbb{F}_3^n \otimes \mathbb{F}_3^n$ by
\[
T(x,y,z) = \begin{cases} 1 & \text{if } x + y + z = 0 \\ 0 & \text{otherwise} \end{cases}
\]
A subset $A$ is a cap set if and only if $T|_{A \times A \times A}$ is \emph{diagonal} (i.e., $T(x,y,z) \neq 0$ implies $x = y = z$). For diagonal tensors, the slice rank equals $|A|$, so
\[
|A| \le \text{slice rank}(T).
\]

\textbf{Step 2: Bounding slice rank by polynomial counting.} The slice rank of $T$ is bounded by $3 \cdot |\mathcal{P}_{<2n/3}|$, where $\mathcal{P}_{<2n/3}$ is the space of polynomials in $\mathbb{F}_3[x_1, \ldots, x_n]$ of degree less than $2n/3$ (with each variable exponent in $\{0,1,2\}$).

\textbf{Step 3: Stirling's approximation.} The number of such monomials is $\sum_{k < 2n/3} \binom{n}{k} 2^k$, which by Stirling's approximation is at most $2.756^n$. Therefore
\[
r_3(n) \le 3 \cdot 2.756^n.
\]

The constant $2.756$ comes from optimizing the entropy bound: it is the exponential rate of $\max_{0 \le t \le 2} 3^{H(t/3)}$ where $H$ is the binary entropy function, evaluated at the critical point.

\subsection{Student-Friendly Explanation: How to Think About the Key Insight}

Here is the conceptual picture, stripped of technicalities:

\textbf{The old way (Fourier analysis).} Think of the cap set condition as saying: ``the set $A$ has no three-term arithmetic progression.'' Fourier analysis converts this into a statement about correlations between $A$ and characters of $\mathbb{F}_3^n$. But characters are \emph{one-dimensional} objects (complex exponentials), and the cap set condition is inherently \emph{three-dimensional} (it involves triples). So Fourier analysis can only give bounds that lose a factor of $3^n$.

\textbf{The new way (slice rank).} Instead of Fourier analysis, use \emph{polynomials}. A polynomial of degree $d$ in $n$ variables over $\mathbb{F}_3$ has at most $d \cdot 3^{n-1}$ zeros (Schwartz-Zippel). If you can find a polynomial that vanishes on all of $A$ except one point, then $|A|$ is bounded by the number of zeros, which is controlled by the degree.

The slice rank method makes this precise: it converts the cap set problem into a question about the rank of a tensor, and then bounds that rank by counting monomials. The key insight is that \emph{low-degree polynomials cannot vanish on too many points}, and the cap set condition forces the existence of such polynomials.

\begin{highlightbox}
\textbf{The punchline:} The slice rank method is a powerful generalization of the polynomial method: instead of a single polynomial vanishing on a set, one uses the rank of a tensor. This upgraded tool can handle three-body constraints (like the cap set condition) that defeat two-body tools like Fourier analysis. It has since been applied to many problems including the sunflower lemma, the Erd\H{o}s-Ginzburg-Ziv theorem, and the problem of finding large subsets of $\mathbb{F}_p^n$ with no solutions to linear equations.
\end{highlightbox}

\subsection{Ellenberg-Gijswijt's Proof in Detail}

\begin{theorem}[Ellenberg-Gijswijt \cite{ellenberggijswijt2017}]
$r_3(n) \le 2.756^n$ for sufficiently large $n$.
\end{theorem}

The proof proceeds in three steps:

\begin{enumerate}
\item \textbf{Tensor formulation}: Represent the cap set condition by a tensor $T: \mathbb{F}_3^n \times \mathbb{F}_3^n \times \mathbb{F}_3^n \to \{0,1\}$ defined by $T(x,y,z) = 1$ if $x + y + z = 0$ and $0$ otherwise. A cap set is a set $A$ such that $T|_{A \times A \times A}$ is diagonal (no off-diagonal entries).

\item \textbf{Slice rank bound}: The slice rank of $T$ is at most $3 \cdot M$, where $M$ is the number of monomials $x_1^{e_1} \cdots x_n^{e_n}$ with $e_i \in \{0,1,2\}$ and $\sum e_i < 2n/3$.

\item \textbf{Stirling's formula}: $M \le 2.756^n$, giving the bound $|A| \le \text{slice rank}(T) \le 3 \cdot 2.756^n$.
\end{enumerate}

\section{Impact}
\begin{itemize}
\item \textbf{Polynomial method 2.0}: Slice rank opened a new chapter in the polynomial method, extending it from geometry (Dvir's Kakeya) to combinatorics.
\item \textbf{Roth's theorem over $\mathbb{F}_3$}: The bound $2.756^n$ is exponentially better than previous methods.
\item \textbf{Connections to TCS}: The slice rank is related to the non-deterministic communication complexity of tensor operations.
\end{itemize}

\section{Exercises}

\begin{exercise}[Basic]
Show that a cap set in $\mathbb{F}_3^n$ has no three points on a line. Compute the maximum cap set size for $n=2$.
\end{exercise}

\begin{exercise}[Intermediate]
Explain how the slice rank of a tensor generalizes the rank of a matrix. Why does slice rank bound cap sets?
\end{exercise}

\section{Timeline}

\begin{tabularx}{\linewidth}{lX}
\toprule
\textbf{Year} & \textbf{Event} \\ \midrule
1953 & Roth proves no 3-AP in dense subsets of $\mathbb{Z}$ \\
1995 & Meshulam proves $r_3(n) \le 2 \cdot 3^n / n$ via Fourier analysis \\
2016 & Croot-Lev-Pach use polynomial method for $\mathbb{Z}_4^n$ \\
2016 & Ellenberg-Gijswijt solve cap set with slice rank \\
2022 & Huh wins Fields Medal (log-concavity in combinatorics) \\
\bottomrule
\end{tabularx}

\section{Citations}

\begin{enumerate}
\item Ellenberg, J., \& Gijswijt, D. (2017). ``On large subsets of $\mathbb{F}_3^n$ with no three-term arithmetic progression.'' Annals of Mathematics, 185(1), 339--343.
\item Tao, T. (2016). ``A symmetric formulation of the Croot-Lev-Pach-Ellenberg-Gijswijt capset bound.'' Blog post.
\item Croot, E., Lev, V., \& Pach, P. P. (2017). ``Progression-free sets in $\mathbb{Z}_4^n$ are exponentially small.'' Annals of Mathematics, 185(1), 331--337.
\item Meshulam, R. (1995). ``On subsets of finite abelian groups with no 3-term arithmetic progressions.'' Journal of Combinatorial Theory, Series A, 71(1), 168--174.
\end{enumerate}
